\exercisesfor{7}

\begin{enumerate}
\def\labelenumi{\arabic{enumi}.}
\item
  Let \(\mathrm{G}\) be a Lie group and \(\phi: \mathrm{U} \to \mathrm{G}\) an exponential mapping such that \(\mathbf{Z}\mathbf{U} \subset \mathbf{U}\) and \(\phi(rx) = \phi(x)^r\) for all \(x \in \mathbf{U}\) and all \(r \in \mathbf{Z}\) . If \(p > 0\) , \(\phi\) is an analytic isomorphism of \(\mathrm{U}\) onto \(\phi(\mathrm{U})\) . (If \(\phi(x) = \phi(y)\) , then \(\phi(p^n x) = \phi(p^n y)\) for all \(n \in \mathbf{N}\) and hence \(x = y\) . Let \(\mathrm{W}\) be an open neighbourhood of \(0\) in \(\mathrm{L}(\mathrm{G})\) such that \(\phi^{-1}\) is analytic on \(\phi(\mathrm{W})\) . For all \(s \in \phi(\mathrm{U})\) , there exist \(n \in \mathbf{N}\) and a neighbourhood \(\mathrm{V}\) of \(s\) in \(\phi(\mathrm{U})\) such that \(t \in \mathrm{V} \Rightarrow t^{p^n} \in \phi(\mathrm{W})\) .)
\item
  Let G be a Lie group, \(\phi: \mathrm{U} \to \mathrm{G}\) an exponential mapping with the properties of Proposition 3, V a neighbourhood of 0 in U such that \(\mathbf{ZV} \subset \mathrm{V}\) and \(\psi: \mathrm{V} \to \mathrm{G}\) a tangent mapping at 0 to \(\phi\) such that \(\psi(nx) = \psi(x)^n\) for all \(n \in \mathrm{V}\) and all \(n \in \mathbf{Z}\) . If \(p > 0\) , then \(\psi = \phi \mid \mathrm{V}\) . (Give L(G) a norm. Let \(x \in \mathrm{V}\) . Then \(p^n x\) tends to 0 as \(n\) tends to \(+\infty\) , hence there exists \(\alpha_n > 0\) such that \(\alpha_n\) tends to 0 and \(\|(\phi^{-1} \circ \psi)(p^n x) - p^n x \| \leqslant \alpha_n \|p^n x\|\) . But \(\psi(p^n x) = \psi(x)p^n\) , whence \(\|(\phi^{-1} \circ \psi)(x) - x \| \leqslant \alpha_n \|x\|\) .)
\item
  Let U be the set of invertible elements of A and \(U' = 1 + m \subset U\) .
\end{enumerate}

\begin{enumerate}
\def\labelenumi{(\alph{enumi})}
\item
  Show that \(\mathbf{U}' \subset \mathbf{U}_f\) . (If \(x = 1 + y\) with \(y \in \mathfrak{m}\) , then \(x^{p^n}\) tends to 1 as \(n\) tends to \(+\infty\) by the binomial theorem.)
\item
  Show that \(U_{f}\) is the set of elements of U whose image in A/m is a root of unity. (Use (a).) Hence recover the fact that \(U = U_{f}\) when K is locally compact (A/m is then finite).
\end{enumerate}

\begin{enumerate}
\def\labelenumi{\arabic{enumi}.}
\setcounter{enumi}{3}
\tightlist
\item
  Let \(n \in \mathbf{N}^*\) , \(p\) be a prime number and \(G\) the set of matrices belonging to \(\mathbf{GL}(n, \mathbf{Z}_p)\) all of whose elements are congruent to 1 modulo \(p\) if \(p \neq 2\) (resp. modulo 4 if \(p = 2\) ). Then \(G\) is an open subgroup of \(\mathbf{GL}(n, \mathbf{Z}_p)\) .
\end{enumerate}

\begin{enumerate}
\def\labelenumi{(\alph{enumi})}
\item
  Show that G has no element of finite order \(\neq 1\) .
\item
  Show that every finite subgroup of \(\mathbf{GL}(n, \mathbf{Z}_p)\) is isomorphic to a subgroup of \(\mathbf{GL}(n, \mathbf{Z} / p\mathbf{Z})\) if \(p \neq 2\) (resp. \(\mathbf{GL}(n, \mathbf{Z} / 4\mathbf{Z})\) if \(p = 2\) ).
\end{enumerate}

¶ 5. (a) Let \(\Gamma\) be a compact subgroup of \(G = \mathbf{GL}(n, \mathbf{Q}_p)\) . Show that there exists a conjugate of \(\Gamma\) which is contained in \(\mathbf{GL}(n, \mathbf{Z}_p)\) . (If \(T\) is a lattice of \(\mathbf{Q}_p^n\) with respect to \(\mathbf{Z}_p\) (Commutative Algebra, Chapter VII, § 4, Definition 1), show that the stabilizer of \(T\) in \(\Gamma\) is open in \(\Gamma\) and hence of finite index and \(\sum_{\gamma \in I} \gamma T\) is a lattice which is stable under \(\Gamma\) .)

\begin{enumerate}
\def\labelenumi{(\alph{enumi})}
\setcounter{enumi}{1}
\item
  Deduce that \(\mathbf{G}_f\) is the union of the conjugates of \(\mathbf{GL}(n,\mathbf{Z}_p)\) .
\item
  Deduce that every finite subgroup \(\Phi\) of \(\mathbf{GL}(n, \mathbf{Q}_p)\) has order a divisor of \(a_n(p)\) , where \(a_n(p)\) is defined by
\end{enumerate}

\[
\begin{array}{l l} a _ {n} (p) = (p ^ {n} - 1) (p ^ {n} - p) \ldots (p ^ {n} - p ^ {n - 1}) & \text {if} p \neq 2 \\ a _ {n} (p) = 2 ^ {n ^ {2}} (2 ^ {n} - 1) (2 ^ {n} - 2) \ldots (2 ^ {n} - 2 ^ {n - 1}) & \text {if} p = 2. \end{array}
\]

(Using (a), it can be assumed that \(\Phi \subset \mathbf{GL}(n, \mathbf{Z}_p)\) . Then use Exercise 4.)

\begin{enumerate}
\def\labelenumi{(\alph{enumi})}
\setcounter{enumi}{3}
\tightlist
\item
  Show that every finite subgroup of \(\mathbf{SL}(n, \mathbf{Q}_p)\) has order a divisor of \(s_n(p)\) , where \(s_n(p) = \frac{a_n(p)}{p - 1}\) for \(p \neq 2\) and \(s_n(2) = \frac{a_n(2)}{2}\) .
\end{enumerate}

¶ 6. In this exercise, l denotes a prime number \(\neq2\) . If a is an integer \(\neq0\) , let \(v_{l}(a)\) denote the l-adic valuation of a, i.e.~the largest integer e such that \(a\equiv0\pmod{l^{e}}\) .

\begin{enumerate}
\def\labelenumi{(\alph{enumi})}
\tightlist
\item
  Let \(m\) be an integer \(\geqslant 1\) . We write
\end{enumerate}

\[
\begin{array}{l l} \varepsilon (l, m) = 0 & \text { if } m \neq 0 (\mathrm{mod} (l - 1)) \\ \varepsilon (l, m) = v _ {l} \left(\frac {m}{l - 1}\right) + 1 & \text { if } m \equiv 0 (\mathrm{mod} (l - 1)). \end{array}
\]

Show that, if \(x\) is an integer prime to \(l\) , then

\[
v _ {l} (x ^ {m} - 1) \geqslant \varepsilon (l, m)
\]

and that there is equality if the image of \(x\) in the cyclic group \((\mathbf{Z} / l^2\mathbf{Z})^*\) is a generator of this group.

\begin{enumerate}
\def\labelenumi{(\alph{enumi})}
\setcounter{enumi}{1}
\tightlist
\item
  Let \(n\) be an integer \(\geqslant 1\) . We write
\end{enumerate}

\[
r (l, n) = \sum_ {m = 1} ^ {n} \varepsilon (l, m).
\]

Show that

\[
r (l, n) = \left[ \frac {n}{l - 1} \right] + \left[ \frac {n}{l (l - 1)} \right] + \left[ \frac {n}{l ^ {2} (l - 1)} \right] + \dots
\]

where the symbol \([\alpha]\) denotes the integral part of the real number \(\alpha\) . Show that, if x is an integer prime to l, then

\[
v _ {l} ((x ^ {n} - 1) (x ^ {n - 1} - 1) \dots (x - 1)) \geqslant r (l, n)
\]

and that there is equality if the image of \(x\) in \((\mathbf{Z} / l^2\mathbf{Z})^*\) is a generator of this group.

\begin{enumerate}
\def\labelenumi{(\alph{enumi})}
\setcounter{enumi}{2}
\item
  In the notation of Exercise 5 (c), show that \(l^{r(l,n)}\) is the greatest power of \(l\) which divides all the \(a_{n}(p)\) , for \(p\) prime \(\neq 1\) (or for all sufficiently large \(p\) , which amounts to the same). (Use (b) applied to \(x = p\) ; then choose \(p\) such that its image in \((\mathbf{Z}/l^2\mathbf{Z})^*\) is a generator of this group, which is possible by the arithmetic progression theorem.†)
\item
  Let \(\bar{\Gamma}\) be a finite subgroup of \(\mathbf{GL}(n, \mathbf{Q})\) and let \(l^e\) be the greatest power of \(l\) dividing the order of \(\Gamma\) . Show that \(e \leqslant r(l, n)\) . (Use Exercise 5 to prove that \(l^e\) divides all the \(a_n(p)\) and then apply (c) above.)
\item
  Conversely, show that there exists a finite \(l\) -subgroup \(\Gamma_{i,n}\) of \(\mathbf{GL}(n,\mathbf{Q})\) whose order is \(l^{d,l,n}\) . (Reduce it to the case where \(n\) is of the form \(l^a (l - 1)\) , with \(a\geqslant 0\) . Decompose \(\mathbf{Q}^{\mathrm{n}}\) as a sum of \(l^a\) copies of \(\mathbf{Q}^{j - 1}\) and use this decomposition to give an operation on \(\mathbf{Q}^{\mathrm{n}}\) of the semi-direct product \(\mathrm{H}_{l,n}\) of the symmetric group \(\mathfrak{S}_{la}\) and the group \((\mathbf{Z} / l\mathbf{Z})^{a}\) . Take \(\Gamma_{l,n}\) to be a Sylow \(l\) -group of \(\mathrm{H}_{l,n}\) .)
\end{enumerate}

Show that, if \(n\) is even, \(\Gamma_{l,n}\) is contained in a conjugate of the symplectic group \(\mathbf{Sp}(n,\mathbf{Z})\) .

\begin{enumerate}
\def\labelenumi{(\alph{enumi})}
\setcounter{enumi}{5}
\tightlist
\item
  *Let \(\Gamma\) be a finite subgroup of \(\mathbf{GL}(n, \mathbf{Q})\) which is an \(l\) -group. Show that \(\Gamma\) is conjugate to a subgroup of \(\Gamma_{l,n}\) . (Show first, using a suitable reduction modulo \(p\) , that the reduction of \(\Gamma\) is a finite subgroup of the reduction of a conjugate of \(\Gamma_{l,n}\) ; then use the character of the representation of \(\Gamma\) on \(\mathbf{Q}^{n}\) .) In particular, every subgroup of \(\mathbf{GL}(n, \mathbf{Q})\) of order \(l^{(l,n)}\) is conjugate to \(\Gamma_{l,n*}\) .
\end{enumerate}

¶ 7. In this exercise, let \(v_{2}(a)\) denote the 2-adic valuation of an integer \(a\) .

\begin{enumerate}
\def\labelenumi{(\alph{enumi})}
\item
  Let \(\Gamma\) be a finite subgroup of \(\mathbf{GL}(n,\mathbf{Q})\) . Show that there exists a positive definite quadratic form with coefficients in \(\mathbf{Z}\) , which is invariant under \(\Gamma\) . Deduce (by the same argument as in Exercises 4 and 5) that, for all sufficiently large \(p\) , \(\Gamma\) is isomorphic to a subgroup of an orthogonal group \(\mathbf{O}(n)\) over the field \(\mathbf{F}_p\) (resp. a subgroup of \(\mathbf{SO}(n)\) if \(\Gamma\) is contained in \(\mathbf{SL}(n,\mathbf{Q})\) ).
\item
  Suppose that \(\Gamma\) is contained in \(\mathbf{SL}(n, \mathbf{Q})\) and let \(2^e\) denote the greatest power of 2 dividing the order of \(\Gamma\) . Show that, if \(n\) is odd, \(2^e\) divides all the integers
\end{enumerate}

\[
b _ {n} (p) = \left(p ^ {n - 1} - 1\right) \left(p ^ {n - 3} - 1\right) \dots \left(p ^ {2} - 1\right)
\]

for sufficiently large prime \(p\) . (Use (a) and Exercise 13 of Algebra, Chapter IX, § 6.)

Show that, if \(n\) is even and \(p\) is a sufficiently large prime, \(2^e\) divides the l.c.m. \(b_n(p)\) of the integers

\[
\begin{array}{l} (p ^ {n} - 1) (p ^ {n - 2} - 1) \dots (p ^ {2} - 1) / (p ^ {n / 2} + 1) \\ (p ^ {n} - 1) (p ^ {n - 2} - 1) \dots (p ^ {2} - 1) / (p ^ {n / 2} - 1). \end{array}
\]

(Same method as for \(n\) odd.)

\begin{enumerate}
\def\labelenumi{(\alph{enumi})}
\setcounter{enumi}{2}
\tightlist
\item
  Under the hypotheses of (b), we write
\end{enumerate}

\[
r (2, n) = n + \left[ \frac {n}{2} \right] + \left[ \frac {n}{4} \right] + \dots = n + v _ {2} (n!).
\]

Let A be an integer \(\geqslant 3\) . Show that \(2^{r(2,n)-1}\) is the greatest power of 2 which divides all the \(b_n(p)\) for \(p \geqslant A\) . (Same method† as in Exercise 6; use the existence of a prime number \(p \geqslant A\) such that \(p \equiv 5 \pmod{8}\) .) Deduce the inequality \(e \leqslant r(2,n) - 1\) .

\begin{enumerate}
\def\labelenumi{(\alph{enumi})}
\setcounter{enumi}{3}
\tightlist
\item
  Conversely, let \(C_{n}\) be the subgroup of \(\mathbf{GL}(n,\mathbf{Z})\) generated by the permutation matrices and the diagonal matrices with coefficients \(\pm1\) . The order of \(C_{n}\) is \(2^{n}n!\) and \(v_{2}(2^{n}n!) = r(2,n)\) . If \(\Gamma_{2,n}\) denotes the intersection of a Sylow 2-group of \(C_{n}\) with \(\mathbf{SL}(n,\mathbf{Z})\) , deduce that the order of \(\Gamma_{2,n}\) is \(2^{n(2,n)-1}\) .
\end{enumerate}

\begin{enumerate}
\def\labelenumi{\arabic{enumi}.}
\setcounter{enumi}{7}
\tightlist
\item
  Let \(n\) be an integer \(\geqslant 1\) . We write
\end{enumerate}

\[
\mathrm{M} (n) = \prod_ {l} 1 ^ {r (1, n)},
\]

where the product is taken over all prime numbers \(l\) and the \(r(l, n)\) are defined as in Exercises 6 and 7.

Then \(\mathbf{M}(1) = 2\) , \(\mathbf{M}(2) = 2^3.3 = 24\) , \(\mathbf{M}(3) = 2^4.3 = 48\) ,

\[
\mathrm{M} (4) = 2 ^ {7}. 3 ^ {2}. 5 = 5 7 6 0.
\]

Deduce from Exercises 6 and 7 that the l.c.m. of the orders of the finite subgroups of \(\mathbf{GL}(n,\mathbf{Q})\) (or \(\mathbf{GL}(n,\mathbf{Z})\) , which amounts to the same) is equal to \(\mathrm{M}(n)\) . Consider the same question for \(\mathbf{SL}(n,\mathbf{Q})\) with \(\mathrm{M}(n)\) replaced by \(\frac{1}{2}\mathrm{M}(n)\) .

\begin{enumerate}
\def\labelenumi{\arabic{enumi}.}
\setcounter{enumi}{8}
\item
  Suppose that \(\mathbf{K}\) is locally compact. Let \(\mathrm{G} = \mathbf{GL}(n,\mathbf{K})\) . Let \(\mathbf{G}_1\) be the set of \(g\in \mathbf{G}\) which leaves stable a lattice of \(\mathbf{K}^n\) with respect to A. Let \(\mathbf{G}_2\) be the set of \(g\in \mathbf{G}\) which generate a relatively compact subgroup in G. Let \(\mathbf{G}_3\) be the set of \(g\in \mathbf{G}\) whose eigenvalues in an algebraic closure of \(\mathbf{K}\) are of absolute value 1. Then \(\mathbf{G}_1 = \mathbf{G}_2 = \mathbf{G}_3 = \mathbf{G}_f\) . (Use the argument of Exercise 5 (a).)
\item
  Suppose that K is locally compact. Let G be a standard group of dimension n over K. Let \(\mu\) be a Haar measure on the additive group \(K^{n}\) . Show that \(\mu|G\) is a left and right Haar measure on G. (Use the fact that G is the inverse limit of the \(G(\alpha_{\lambda})\) and Integration, Chapter VII, §1, Proposition 7.)
\end{enumerate}

